\documentclass[]{../../../../tex/classes/styledArticle}
\usepackage{../../../../tex/packages/hebrewSupport}
\usepackage{../../../../tex/packages/mathShortcuts}
\usepackage{../../../../tex/packages/theoremsSupport}

\author{שחר פרץ}
\title{אלגוריתמים $\sim$ \textit{הרצאה 9}}
\begin{document}
	\maketitle
	\textbf{מרצה: }עמית ויינשטין
	
	\subsection*{תכנות לינארי}
	זה פטיש חזק שמסוגל לפתור את רוב הבעיות בקורס עד עכשיו (גם אם לא בסיבוכיות טובה). נראה אלגוריתם שברמה הפרקטית רץ טוב, למרות שאין דרך לפרמל את זה (אפילו לא בתוכלת). ניתן לפתור את הבעיות האלו בזמן פולינומיאלי כי אם בפועל הפתרון הזה פחות יעיל. 
	
	בגדול – יש $n$ משתנים $x_1 \dots x_n$, ו־$m$ א''שים ($\le, =, \ge$). כמו כן תהיה לנו פונקציית מטרה שתנסה למזער או למקסם צירוף לינארי של המשתנים, תחת האילוצים (הא''שים). 
	
	לדוגמה, נרצה למצוא את $\max x_1 + 2x_2$ כך ש־$4x_1 \le 3 \land 2x_1 - 2x_2 \ge 0$ (במקרה זה $m = n = 2$, שני משתנים $x_1, x_2$ ושני א''שים). 
	
	\defi{פתרון נקרא \textit{פיזיבילי} אם הוא עומד באילוצים. נאמר שנקודה היא \textit{נקודה פיזיבילית} אם הפוליגון שנוצר ע''י האילוצים לא דיפרנציאבילי בנקודה זו. }
	
	ייתכנו כל מיני מצבים: 
	\begin{itemize}
		\item אין פתרון
		\item יש פתרון אופטימלי (לאו דווקא יחיד, הוא יכול להיות פאה של האובייקט)
		\item הבעיה לא חסומה (יש פתרון פיזיבילי טוב כרצוננו). 
	\end{itemize}
	
	
	\subsubsection*{זרימות}
	נבחין שניתן להציג כל מיני בעיות שאנו מכירים כבעיות תכנות לינארי. לדוגמה, עבור זרימות (בהינתן $G =(V, E)$ מכוון ו־$c \co E \to \R^{+}$) לכל קשת $(u, v) \in E$ נגדיר משתנה $x_{u, v}$, תחת האילוצים: 
	\begin{itemize}
		\item חוק שימור החומר: $\forall u \in V \setminus \{s, t\} \co \sum_{v \in V}x_{u, v} = 0$
		\item קיום שמורת הקיבול: $\forall (u, v) \in E \co x_{u, v} \le c(u, v)$
		\item אנטי־סימטריה: $\forall u, v \in E \co x_{u, v} = -x_{v, u}$
	\end{itemize}
	ופונקציית המטרה תהיה ערך הזרימה, שהוא $f(x) = \sum_{v \in \adj[v]}x_{s, v}$. 
	
	זמן הריצה שנקבל מהפעלת תכנות לינארי יהיה יותר גרוע (במיוחד ברשתות 01), אבל לפחות נוכל להבטיח פתרון פולינומיאלי. 
	
	נוכל ''לשדרג`` את הבעיה ולהוסיף $p \co E \to \R^{+}$ מחיר להזרמת חומר על כל עלות. ואז נקבל עוד פונקציה $\sum_{u, v \in E}p_{u, v}x_{u, v}$ שנרצה למזער, בעוד את הפונקציה השנייה נרצה למקסם. איך נעשה זאת? כדי למצוא זרימה אופטימלית במחיר מינימלי, נרצה קודם למצוא זרימה אופטימלית בערך $T$, להוסיף אילוץ לינארי שזו ערך הזרימה, ולהחליף את פונקציית המטרה במחיר עם מינימלי. 
	
	\subsubsection*{מק''בים}
	בהינתן $G = (V, E)$ מכוון, $w \co E \to \R$ משקל, ו־$s, t \in V$, כיצד נמצא מסלול קל ביותר $s \squig t$. בשביל הבעיה הזו, נרצה למצוא את $\dd[v]$ ($\dg(s, v)$) לכל קודקוד. נגדיר משתנה $x_v$  לכל $ v\in V$. האילוצים יהיו שלכל קשת $(u, v)$, מתקיים $x_u \le x_v + w(u, v)$ וכן $x_s = 0$ (אחרת זה יהיה unbounded). 
	
	לגבי פונקציית המטרה, הבעיה תהיה unbounded אם נרצה למצוא מינימום של $x_t$. נרצה למצוא את המקסימום של $x_t$. 
	
	למה זה עובד? כל מסלול $s = v_0 \squig v_k = t$ מתקיים $x_t \le w(v_{k - 1}, v_k)$ ובאינדוקציה $x_t \le \sum_{i = 0}^{k -1} w(v_i, v_{i + 1})$. בשפה טבעית, כל מסלול חוסם מלמעלה את $x_t$. נוכל לקחת פשוט את המקסימום ונקבל את המק''ב. 
	
	\subsubsection*{הצורה הסטנדרטית}
	סימפלקס הוא הרחבה של הרעיון של נקודה, משולש פרמידה וטטהרדרון ומספר ממדים. 
	
	בשביל אלגו' Simplex, ראשית נגדיר את הצורה הסטנדרטית לבעיית תכנות לינארי: 
	\begin{itemize}
		\item פונקציית ממטרה ממקסמת
		\item כיוון א''ש $\le$ (כלומר $\sum_{j = 1}^{n} a_{ij}x_j \le b_j$ ל־$i = 1 \dots m$) 
		\item נדרוש שכל המשתנים אי־שליליים (כלומר $x_j \ge 0$)
	\end{itemize}
	כל בעיית LP ניתן לכתוב בצורה הסטנדרטית: 
	\begin{itemize}
		\item כדי להגיע לפונקציית מטרה ממקסמת, אפשר להפוך את הכיוון של המקדמים
		\item אין בעיה להפוך אי שיוונות מהצורה $\ge$, ושיוויונות אפשר להמיר לשני אי־שיוונות
		\item כל משתנה אי־שלילי אפשר להציג בתור חיסור של שני משתנים: $x \mapsto x' - x''$ כאשר $x', x''$ אי־שליליים. 
	\end{itemize}
	סה''כ $n$ משתנים ו־$m$ אילוצים הפכו להיות לכל היותר $2n$ משתנים ולכל היותר $2m$ אילוצים. 
	
	השימוש בצורה הסטנדרטית מאפשר לנו להמיר את הבעיה לייצוג מטריציוני: מטריצה $A = \{a_{ij}\}_{m \times n}$ של מקדמי האילוצים, וקטור $\{b_i\}_{m} = b$ של מקדמים חופשים, ו־$\{c_j\}_{n} = c$ של מקדמי הפתרון, ווקטור $x = \{x_i\}_{i = 1}^{n}$ שמייצג את הפתרון. ואז נדרוש ש־$Ax \le b$ ונמקסם את $c^Tx$
	
	\subsubsection*{אלגוריתם Simplex (1947)}
	צורת Slack -- נעבוד עם המרווחים (Slack) של כל אחד מהאילוצים. 
	
	נגדיר משתנה חדש $x_{n + i}$ שיאמר לנו עד כמה הא''ש לא הדוק: 
	\[ \sum_{j = 1}^{n} a_{ij} x_j \le b_i \iff 0 \le x_{n + i} = b_i - \sum_{j =1}^{n}a_{ij}x_j \]
	בניסוח מטריציוני, הגדרנו וקטור $s$ המוגדר לפי $s = b - Ax$ ועתה נדרוש ש־$s \ge 0$ כדי לוודא שאנו עומדים בתנאי הא''שים (ואפשר להתייחס ל־$s$ כאל שרשור ל־$x$). 
	
	ל־$m$ האלמנטים של $s$ נקרא \textit{משתנים בסיסיים} (הם $x_{n + i}$ כאשר $i = 1 \dots m$). 
	
	כמובן שעדיין נרצה למקסם את $c^{T}x$, פשוט תחת האילוץ ש־$s,  x \ge 0$. 
	
	נוכל לבצע החלפה בין משתנים בסיסיים למשתנים לא בסיסיים ע''י החלפת משתנים. 
	
	\textit{הערה: }יש כאן שורה של דברים שממש קשה להסביר בקובץ latex אז כאן תנסו להשלים את האלגוריתם מההרצאה. בכל מקרה, זה מה שהמרצה כתב כדי לנסות להבהיר את זה בצורה כתובה. 
	
	בצעד Pivot נעבור בין צורוך Slack שבתקווה משפרות את ערך פונקציית המטרה. בוחרים משתנה לא בסיסי עם מקדם חיובי (בשורה למעלה) (שכן מקדם שלילי רק יכול לגרוע) ונחליף אותו עם משתנה בסיסי המביא למינימום את: 
	\[ \min_{i \in B} \frac{b_i}{a_{i1}} \]
	כדי לוודא שהפתרון הבא גם הוא פיזיבלי (בהנחה שהתחלנו מפתרון פיזיבילי)
	
	\textit{הבחנה: }אם אנחנו בקודקוד, אז יש לנו $n$ שיוויונות הדוקים (כי $n$ הממד שלנו). 
	
	נוכל לעבור בין צורות Slack שונות (למעבר כזה נקרא צעד pivot), ואם נציב אפסים, צורת ה־Slack הזו תשרה פתרון אפשרי. 
	
	עולה השאלה – כמה נקודות יש? נוכל במקום זאת לשאול כמה קומבינציות של $n$ משוואות הדוקות ל־$m$ האילוצים שלנו יש. כמה כאלו יש? $\binom{m}{n}$, וזה המון...  אבל לא כל החלפת משנתה היא רלוונטית: ידוע שהמשתנים הבסיסיים אי־שליליים. לכן, לא רלוונטי יהיה לבצע pivot על משתנה עם מקדם שלילי (כי במשוואה יהיה מקדם חופשי חיובי, ואז נקטין את המשתנה החופשי בשורה של המשתנים הבסיסיים). אינטואטיבית, זה מה שהופך את האלגוריתם יעיל בפרקטיקה. 
	
	ערך הפתרון בהכרח לא קטן, אבל יכול להשאר זהה. על מנת להמנע מלולאה אינסופית, מספיק לבחור משתנה לא בסיסי עם אינדקס מינימלי. 
	
	אז כיצד בפועל עובד סימפלקס? נעבור לצורת slack, ונבצע מעברי pivot עד שמגלים שהבעיה לא חסומה או שמגיעים לפתרון אופטימלי. מעבר pivot אורך $\oc(nm)$, ולכן במקרה הגרוע הסיבוכיות תהיה $\oc\cl{\binom{n + m}{n}nm}$ למרות שבפועל האלגוריתם מהיר. 
	
	עד עתה הנחנו שהצורה ההתחלתית היא פיזיבלית. אבל איך יודעים אם הצורה ההתחלתית אכן פיזבלית? נוסיף משתנה $x_0$, ונבדוק ש־$\forall 1 \le i \le m \co \sum_{j = 1}^{n}a_{ij}x_j - x_0 \le b_i$ (כאשר לכל $j \in 0 \dots n$ מתקיים $x_j \ge 0$), ואז לחפש את המקסימום של $x_0$. 
	
	\theo{ערך הפתרון בבעיית העזר לעיל הוא $0$ אמ''מ הבעיה המקורית פיזיבילית. }
	
	\prooftrivial
	
	מהסתבר שצעד pivot אחד על בעיית העזר מוביל את בעיית העזר לפתרון הפיזיבילי (ואז נוכל לוודא שערך המקסימום של $x_0$ הוא אכן $0$). 
	
	השאלה העיקרית שנותר לנו לעלות עליה הוא להראות שסימפלקס אכן מוצא לבסוף את הפתרון הפיזבילי. לשם כך נגדיר את הבעיה הדואלית – נחליף את האילוצים במשתני המטרה, ונראה שזה עובד. 
	
	\begin{itemize}
		\item בבעיה המקורית נרצה למצוא את $\max \sum_{j = 1}^{n}c_j x_j$, בדואלית נרצה למצוא את $\min \sum_{i = 1}^{m}b_i y_i$. 
		\item בבעיה המקורית האילוצים שלנו הם $\sum_{j = 1}^{n}a_{ij}x_j \le b_i$. בבעיה הדואלית נרצה לפתור את $\sum_{i = 1}^{m}a_{ij}y_i \ge c_j$
		\item בבעיה המקורית נדרוש ש־$x_i \ge 0$. בבעיה הדואלית נדרוש ש־$y_i \ge 0$. 
	\end{itemize}
	
	\textbf{איטואיציה – זיווג שברי: }לכל קשת נגדיר משתנה $0 \le x_e \le 1$. לכל קודקוד $v$ נאמר ש־$\sum_{e \ni v}x_e \le 1$ ונמקסם את $\max \sum x_e$ (זו הרחבה של בעיית הזיווג למצב בו מאשרים לקחת ''חצי קשת``. נרצה למצוא אוסף מקסימלי של קשתות שהן זרות בקודקודים). בבעיה הדואלית נרצה למצוא אוסף מקסימלי של קודקודים שנוגעים בכל הקשתות (קוראים לזה vertex cover). 
	
	
	נסמן ש־$\bar x_1 \dots \bar x_n$ פתרון אופטימלי לפרימרית ו־$\bar y_1 \dots \bar y_m$ פתרון אופטימלי לדואלית. 
	\begin{Theorem}[משפט הדואליות החלש]
		האופטימום של הפרימרי קטן או שווה מהאופטימום של הדואלי
	\end{Theorem}
	\begin{proof}
		אכן, ראשית מהאילוים של הדואלית ולאחר מכן מהאילוצים של הפרימרית: 
		\[ \sum_{j = 1}^{n} c_j \bar x_j \le \sum^{n}_{j = 1} \bar x_j \cl{\sum^{m}_{i = 1} a_{ij}\bar y_i} = \sum_{i = 1}^{m}\cl{\sum_{j = 1}^{n}a_{ij}\bar x_j}\bar y_i \le \sum_{i = 1}^{m}b_i\bar y_i  \]
	\end{proof}
	\cola{בפרט, כל פתרון פיזיבלי בבעיה הדואלית גדול או שווה לכל פתרון פיזיבלי בבעיה הפרימרית. }
	
	\begin{Theorem}[משפט הדואליות החלש]
		מתקיים ש־$\mathrm{opt}_p = \mathrm{opt}_D$, האופטימום של הדואלית שווה לאופטימום של הפרימרית ($\sum c_j \bar x_j = \sum b_i \bar y_i$). 
	\end{Theorem}
	נוכיח את המשפט ביחד עם ההוכחה ש־simplex עוצר על פתרון אופטימלי, ונראה מה הפתרון של הבעיה הדואלית. 
	\begin{proof}
		נניח שסימפלקס עצר על הפתרון: 
		\[ \max v + \sum_{j \in N}c_j' x_j \]
		כאשר $N$ קבוצת המשתנים הלא־בסיסיים ו־$B$ קבוצת הבסיסיים. עוד נסמן שלכל $i \in B$: 
		\[ x_j = b_i' - \sum_{j \in N}a_{ij}' x_j \]
		משום שסימפלקס עצר בבירור $c_j' \le 0$. נגידר פתרון (אופטימלי) עבור הבעיה הדואלית: 
		\[ \bar y_i = \begin{cases}
			- c_{i + n}' & n + i \in N \\
			0 & n + i \in B
		\end{cases} \]
		נראה שערך הפתרון שווה ל־$v$ (ערך הפתרון של הפרימרית). כל הצבה בצורת slack אחת תביא את אותו הערך בצורה אחרת, כי כל צורות ה־slack שקולות. לכל הצבה של $x_1 \dots x_n$ קיימת קבוצת ערכים $x_{n + 1} \dots x_{n + m}$ כך ש־: [נגדיר $c_j' = 0$ עבור $j \in B$]
		\[ \sum_{j = 1}^{n} c_j x_j = V + \sum_{j \in N} c_j' x_j = V + \sum_{j = 1}^{n + m}c_j' x_j = V + \sum_{j = 1}^{n}c_j'x_j + \underbrace{\sum_{i = 1}^{m}c_{n + i}'x_{n + i}'}_{(\mathrm I)} = * \]
		נבחין ש־: 
		\[ (\mathrm I) = \sum _{i = 1}^{m}- \bar y_i x_{n + i} = \sum_{i = 1}^{m} - \bar y_i\cl{b_i - \sum_{j = 1}^{n}a_{ij} x_j} \]
		נציב חזרה: 
		\[ * = V + \sum_{j = 1}^{n}c_j' x_j + \sum_{i = 1}^{m}- \bar y_i \cl{b_i - \sum_{j = 1}^{n}a_{ij} x_j} = \cl{v - \sum_{i = 1}^{m}\bar y_i b_i} + \sum_{j = 1}^{n}\cl{c_j'+\sum_{i = 1}^{m}a_{ij}\bar y_i}x_j \]
		סה''כ מטרנזטיביות $\sum c_j x_j = V - \sum \bar y_i b_i + \sum\cl{\mathrm{blatttt}}x_j$. לכן לאחר שנציב $x_j = 0$ בכל ה־$j$־ים נקבל ש־: 
		\[ V = \sum_{i = 1}^{m}b_i \bar y_i  \]
		ולכן ערכי $\bar y_i$ מביאים פתרון אופטימלי לדואלי. 
		
		בבירור מהגדרתם ה־$\bar y_i$־ים אי־שליליים. 
		
		עתה נדרש להוכיח ש־$\bar y_i$ אכן פיזבולו. עבור כל $j$ ספציפי, נציב $x_j= 1$ ובשאר $0$. נציב במשוואה ממקודם ונקבל: (כי $c'_j \le 0$)
		\[ c_j = \sum_{i = 1}^{n}a_{ij} \bar y_i \le \sum_{i = 1}^{m}a_{ij} \bar y_i  \]
		וזהו
	\end{proof}
	
	
	
	\ndoc
\end{document}